\section*{ID8404: What Survives: Cascade + Cutoff}
\paragraph{Metadata.}
PiID: 6257591633303 | RTEID: RTE:P32634.5983:T3.5063 | UUID: 090aa3d6-5813-5a55-8c90-e2292c50249e

\paragraph{Canonical Statement.}
In both cases: \textbackslash{}[ \textbackslash{}text\{cascade\} + \textbackslash{}text\{cutoff\} = repeated geometry \textbackslash{}]

\paragraph{Canonical Equation.}
\[
\text{cascade} + \text{cutoff} = repeated geometry
\]

\paragraph{Dictionary.}
Relation: cascade + cutoff = repeated geometry

\paragraph{Encyclopedia.}
What Survives: Cascade + Cutoff is an algebraic definition, written cascade + cutoff = repeated geometry. It is an algebraic definition expressing the quantity directly in terms of the variables on its right-hand side. It is built from r, p, a, g, o, m, defining cascade + cutoff. Within the Trawin Zero Point registry it serves as one element of the framework's mathematical narrative.
